\section{Impact Statement}\label{impact statement}


Our discoveries provide strong empirical evidence for validating the potential to simulate the trust behavior of humans with LLM agents, and pave the way for simulating more complex human interactions and  social institutions where trust is an essential component. 

Simulation is a widely adopted approach in multiple disciplines such as sociology, psychology and economics~\citep{ziems2023can}. However, conventional simulation methods are strongly limited by the expressiveness of utility functions~\citep{ellsberg1961risk,machina1987choice}. Our discoveries have illustrated the great promise of leveraging LLM agents as the simulation tools for human behavior, and have broad implications in social science, such as validating hypotheses about the causes of social phenomena~\citep{easley2010networks} and predicting the effects of policy changes~\citep{kleinberg2018human}.



Another direction of applications for human simulation is to use LLMs as role-playing agents, which can greatly benefit humans~\citep{yang2024social,chen2024persona,Shanahan2023role,ma2024students}. For example, \cite{shaikh2024rehearsal} proposed to let individuals exercise their conflict-resolution skills by interacting with a simulated interlocutor. \cite{yue2024mathvc} developed a virtual classroom platform with simulated students, with whom a human student can practice his or her mathematical modeling skills by discussing and collaboratively solving math problems. 




However, this paper also shows that some LLMs, especially the ones with a relatively small scale of parameters, are still deficient in accurately simulating human trust behavior, suggesting the potential to largely improve their behavioral alignment with humans. In addition, our paper also demonstrates the biases of LLM agents' trust behavior towards specific genders and races, which sheds light on the potential risks in human behavior simulation and calls for more future research to mitigate them. 





